\documentclass[11pt,a4paper]{article}

\usepackage[T1]{fontenc}
\usepackage[utf8]{inputenc}
\usepackage[english]{babel}
\usepackage{lmodern}
\usepackage{microtype}
\usepackage[a4paper,margin=2.4cm]{geometry}
\usepackage{enumitem}
\usepackage[autostyle]{csquotes}
\usepackage[hidelinks]{hyperref}
\usepackage{amsmath}
\usepackage[backend=biber,style=authoryear,maxcitenames=2,maxbibnames=12]{biblatex}
\addbibresource{../bibliographie/references.bib}

\setlength{\parindent}{0pt}
\setlength{\parskip}{0.65em}
\setlist[itemize]{leftmargin=1.5em,itemsep=0.35em,topsep=0.35em}
\setlist[enumerate]{leftmargin=1.8em,itemsep=0.5em,topsep=0.4em}

\newcommand{\projectversion}{3}
\newcommand{\projectdate}{28 July 2026}

\title{The Emergence of the Interesting\\
\large Birth of Ideas, Understanding, and Singularity of Forms}
\author{François Pachet}
\date{Thesis project -- version \projectversion\\\small \projectdate}

\begin{document}

\maketitle

\begin{abstract}
This project aims to establish the interesting as a philosophical object. It
starts from an ordinary yet elusive experience: an idea emerges, a form becomes
intelligible, a passage holds our attention and calls for renewed engagement.
The central hypothesis is that what is interesting, for a situated subject and
relative to a given horizon, triggers and sustains a construction that provides
a new grip on an object, a problem, or a space of possibilities. The interesting
is thus neither an intrinsic quality of a form nor a private preference, but a
dispositional, historical, and temporal relation. The thesis will reconstruct
its genealogy, identify its properties, and test the concept in five
laboratories: music, artificial intelligence, creativity, philosophy, and
literature. Its reflexive dimension will be decisive: the concept of the
interesting can itself become interesting and transform the subject attempting
to grasp it.
\end{abstract}

\section{Origin, Question, and Aim}

This work emerges from a constructivist and empirical practice of artificial
intelligence research. I have spent forty years building systems that generate,
transform, or accompany music and text. The doctoral project is not intended to
summarise this work or retrospectively impose an artificial philosophical unity
on it. It starts from a more demanding question: \emph{what am I doing} when I
build a system, recognise a solution, reject an output that is correct but
uninteresting, or stop before a form that seems to impose itself?

The guiding question is:

\begin{quote}
Under what conditions does something appear interesting to a subject, sustain
the subject's activity, and then lose or transform this power?
\end{quote}

A song, a proof, a sentence, or an idea may hold attention disproportionately
and provoke comparisons, repetitions, and variations. This effect cannot be
reduced to raw novelty, because randomness surprises without necessarily being
interesting; nor to beauty, because a familiar form may remain beautiful while
becoming predictable; nor to popularity; nor to subjective preference
understood as a datum without a history.

The aim is threefold. The thesis must first establish the philosophical place of
the interesting among competing notions. It must then construct a concept
sufficiently determinate to produce distinctions and counterexamples. Finally,
it must show what the concept makes intelligible in actual practices, without
automatically turning a scientific result or a personal experience into
philosophical proof.

\part{Establishing the Interesting as a Philosophical Problem}

\section{A Blind Spot Rather Than an Absence}

Philosophy has not simply ignored the interesting. It has addressed several of
its dimensions under other names: access to the intelligible in Plato and
Aristotle
\parencite{platon2016republique,platon1995theetete,aristote1991metaphysique,aristote2005seconds};
intuition and evidence in Descartes and Spinoza
\parencite{descartes2002regles,descartes1992meditations,spinoza1966ethique};
judgement in Kant \parencite{kant2001raison,kant1993juger}; invention and the
reconfiguration of problems in Peirce, Poincaré, Hadamard, Bachelard, and Valéry
\parencite{peirce1955writings,poincare1908science,hadamard1959invention,bachelard1999formation,valery1973cahiers};
memory, duration, and the temporal constitution of meaning in Bergson, Husserl,
and Merleau-Ponty
\parencite{bergson2013evolution,bergson2013pensee,husserl1970experience,husserl1957logique,merleau1945phenomenologie};
aspect change in Wittgenstein \parencite{wittgenstein2004recherches}; and
symbolisation, difference, and repetition in Goodman and Deleuze
\parencite{goodman1976languages,deleuze1968difference}.

The proposed diagnosis is therefore not one of lexical absence but of
\emph{functional dispersion}. Truth evaluates the validity of a proposition;
beauty, a particular regime of aesthetic judgement; pleasure, a tone of
experience; surprise, a departure from expectation; curiosity, an orientation
towards missing information; relevance, a relation to a question or goal; and
preference, an ordering of choices. Each concept captures a genuine dimension.
None explains by itself why a correct, intelligible, true, or pleasant form
triggers an activity in a subject that transforms their representations and
possibilities for action.

The interesting may name this level of selection and transformation. It does
not replace neighbouring notions: it allows us to examine their articulations
and divergences. One can like a piece of music without finding it interesting,
be surprised without gaining a new grip, recognise a proposition as true
without judging it worth pursuing, or be mobilised by an object that is painful
and morally objectionable.

\section{State of the Art and Reading Programme}

A historical genealogy is being assembled. It will examine, in particular, the
aesthetics of the interesting in Christian Garve; its opposition to beauty and
relation to modernity in Friedrich Schlegel; Schopenhauer's essay
\emph{Ueber das Interessante}; the relation between boredom, variation, and
\emph{det Interessante} in Kierkegaard; and the articulation of the interesting
and truth in Whitehead. The historical lexicon provided by
\textcite{nannini2018interesting} offers an initial guide, but it cannot replace
a return to the primary texts.

In the twentieth and twenty-first centuries, a small number of works have
directly attempted to conceptualise the interesting. Stace, Kolnai, and Ngai
have analysed its aesthetic value, apparent weakness, and discursive circulation
\parencite{stace1944interestingness,kolnai1964concept,ngai2008merely,ngai2012categories}.
\textcite{epstein2009interesting} stresses its transdisciplinary structure, and
\textcite{grimm2011interesting} its epistemic value in selecting questions.
These works prevent any premature claim to have discovered an untouched
continent. They suggest instead a weak canonisation: the interesting is present,
but without a common architectonic comparable to those of truth or beauty.

This state of the art will remain methodologically stratified. A verified
bibliographic record, an identified passage, a work actually read, and an
interpretation integrated into the argument represent different levels of
acquisition. Sources already read may support the text; the others form an
explicit programme of reading and verification. This distinction prevents a
bibliographic suggestion from being treated as an established authority.

\section{The Interesting and the True}

The interesting cannot become a relativistic substitute for truth. A
recollection of a discussion with Heinz Wismann states the difficulty bluntly:

\begin{quote}
Deleuze et consorts étaient stimulants mais faux. Mais que la philosophie se
devait d'être vraie meme si ennuyeuse
\end{quote}

The quotation distinguishes two values. What is interesting or stimulating
recommends that a construction or inquiry be pursued; truth warrants retaining
a proposition as valid. A false proposition may be heuristically productive
without this productivity making it true. Conversely, a true proposition may
be boring without ceasing to be true.

This distinction does not make the interesting philosophically superfluous.
Among possible truths, many are trivial or irrelevant to a given inquiry. Truth
alone does not select the problems to which attention should be devoted. The
interesting operates prospectively in the choice of questions, hypotheses, and
constructions; truth operates in their validation. Philosophy may begin from a
stimulating error, but it cannot turn the fertility of that error into
permission to remain in falsehood.

\section{The Hostility of Formalisation}

Science has powerful instruments for making a problem tractable: operational
definition, decomposition, control of variables, repeatability, normalisation,
averaging, benchmarking, and optimisation. Yet these instruments also tend
silently to replace the original problem with a more stable substitute.
``Producing something interesting'' then becomes producing something
syntactically correct, semantically coherent, probable in a corpus, surprising,
or compliant with constraints.

In this sense there is a \emph{structural hostility of formalisation towards the
interesting}. This is neither a hostile intention on the part of researchers nor
an anti-scientific argument. The difficulty follows from the relational,
historical, and transformative nature of the phenomenon: making it measurable
may neutralise the relation, history, and transformation that constitute it.
The warning against premature mathematisation
\parencite{russell1995rationality,russell1995awardlecture} is a guiding case.

Each measure must therefore be treated as a local hypothesis. A successful
formalisation must not merely make the problem solvable; it must state what it
preserves from the question and what it abandons. The thesis will thus seek a
reflexive formalisation capable of retaining its residuals instead of declaring
them unimportant.

\part{Constructing the Concept}

\section{A Relational and Constructive Definition}

The current candidate for the core of the thesis is:

\begin{quote}
What is interesting is that which, for a situated subject and relative to a
given horizon, triggers and sustains a construction that provides a new grip on
an object, a problem, or a space of possibilities.
\end{quote}

The minimal relation may be written
\[
  I(F,S\mid H,t),
\]
where \(F\) is the form encountered, \(S\) the subject, \(H\) the social,
cultural, and technical horizon, and \(t\) the moment of encounter. The
collective horizon is not added externally to the two poles: it shapes the
subject's expectations and abilities as well as the repertoire of available
forms. It becomes an autonomous third term when the inquiry concerns the social
circulation or stabilisation of interest.

The term \emph{construction} refers here to an activity, sometimes implicit, and
not necessarily to the fabrication of a material object. It may be:

\begin{itemize}
  \item perceptual and conceptual, when the subject discriminates invariants,
        forms a category, or changes a partition;
  \item explanatory, when the subject arranges relations into a model or
        reconstructs the problem of which the object appears as a solution;
  \item operational, when the subject produces a comparison, prediction,
        intervention, algorithm, or artefact.
\end{itemize}

An arbitrary construction is not enough. The new grip must be manifestable
through a distinction, transfer, anticipation, controlled variation,
intervention, or artefact whose successes and failures can be examined. The
interesting nevertheless guarantees neither the truth of this grip nor its
moral value: it describes its power to initiate and transform.

\section{Candidate Properties}

The constructive definition does not yet constitute an axiomatics. It identifies
a cluster of candidate properties, neither all independent nor necessary in
every case. Their function is to produce variations, objections, and
comparisons.

\subsection{Relational and Dispositional Status}

\begin{enumerate}
  \item \textbf{Relational indexicality.} The interesting is neither contained
        in the form alone nor reducible to a private state. It varies with
        \(F\), \(S\), \(H\), and \(t\).
  \item \textbf{Dispositionality.} A form may be capable of triggering a
        construction without actually doing so, for lack of attention, time, or
        competence. The disposition is actualised only under certain conditions
        of encounter.
  \item \textbf{Path dependence.} The order of experiences matters. A form
        encountered too early may remain opaque; after a new grip has been
        acquired, it may become interesting.
\end{enumerate}

\subsection{Dynamics of Construction}

\begin{enumerate}[resume]
  \item \textbf{Assimilable gap.} There must be enough difference to call for a
        construction and enough grip for it to begin. The familiar tends
        towards triviality; the absolutely foreign towards noise or
        abandonment.
  \item \textbf{Representational productivity.} An interesting form generates
        more comparisons, hypotheses, questions, scenarios, or distinctions
        than it explicitly states.
  \item \textbf{Structured incompleteness.} A gap is interesting only if it
        calls for a directed continuation. The object provides clues for
        constructing what is missing or reformulating the question.
  \item \textbf{Transformation of the space of possibilities.} A strong
        construction does not merely add a representation: it changes the
        questions that can be asked, the distinctions available, or the actions
        that can be conceived.
  \item \textbf{Non-monotonicity.} More information, surprise, or attention does
        not imply more interest. An explanation may renew the construction or
        close too quickly the indeterminacy that sustained it.
  \item \textbf{Self-extinction or renewal.} Some relations are consumable:
        resolution removes their interest. Others are generative: each new grip
        reveals a resistance or a further question.
\end{enumerate}

\subsection{Value, Transfer, and Circulation}

\begin{enumerate}[resume]
  \item \textbf{Relative affective neutrality.} The interesting may be
        pleasant, painful, disturbing, or repellent. It does not in itself
        constitute moral or aesthetic approval.
  \item \textbf{Cross-domain resonance.} A construction gains scope when it
        becomes a schema, model, or controlled metaphor for other domains.
  \item \textbf{Generative compression.} Some forms retrospectively condense
        several phenomena while prospectively opening new derivations.
  \item \textbf{Mediated contagion.} Interest is not transmitted as a property
        attached to an object. Yet one may prepare the conditions for its
        appearance in others by establishing an expectation or providing
        grips.
\end{enumerate}

\subsection{Reflexive Property}

\begin{enumerate}[resume]
  \item \textbf{Self-application.} The concept of the interesting may itself
        trigger the comparisons, distinctions, models, and revisions it
        describes. ``The interesting is interesting'' claims that the concept
        may, under certain conditions, enter its own domain of application.
\end{enumerate}

\section{Temporality, Learning, and the Almost Learnable}

Interest moves with learning. A pattern that is initially opaque may become
accessible, produce a gain in understanding, and then lose its force once
acquired. Theories of compression progress and intrinsic motivation formalise
one version of this dynamic
\parencite{schmidhuber1997interesting,schmidhuber2009compression,oudeyer2007intrinsic,oudeyer2007typology}.
If \(E_t(o,S)\) is the error of the model held by subject \(S\) when facing
object \(o\), local progress may be written
\[
  P_t(o,S)=E_{t-\Delta}(o,S)-E_t(o,S).
\]
Situations already mastered and situations that cannot be learned both have
progress close to zero. Progress niches are those in which improvement is
possible.

This hypothesis connects with, but is not reducible to, the intermediate zone
between boredom and anxiety. With fixed skills, increasing difficulty yields
the sequence
\[
  \text{boredom}\longrightarrow\text{interesting/flow}
  \longrightarrow\text{anxiety}.
\]
For a fixed task, learning may reverse the path. The flow zone refers here to a
structural adjustment between competence and difficulty, not only to the
psychological state of absorption. Hesitation or rupture can remain interesting
as long as the subject retains a grip and can reorganise expectations.

Language shifts this frontier. A word, concept, or notation compresses
distinctions and sometimes operations into a reusable unit. What once required
work becomes a grip for a new construction. The zone of the almost learnable is
therefore historically mobile and symbolically mediated
\parencite{vygotsky1978mind,lenat1991thresholds,lenat2008turtle,koestler1989act,steels2005coordinating,steels2015talkingheads}.
Naming alone is insufficient: a term is a grip only if it can be unfolded into
distinctions, comparisons, predictions, operations, or controlled examples.

\section{Recursion: ``The Interesting Is Interesting''}

The relation first possesses a diachronic recursion:
\[
  I(F,S\mid H,t)\longrightarrow
  \text{construction}\longrightarrow S'
  \longrightarrow I(F,S'\mid H,t+1).
\]
Construction modifies the subject and therefore the conditions of later
evaluations. The object is no longer encountered by the same subject in the
relevant sense: their categories, expectations, and space of possibilities
have changed.

It also possesses a metaconceptual recursion. It may become interesting to
understand why something is interesting, and then to understand what this first
explanation still leaves unexplained. The inquiry then takes as its new object
the relation that triggered it. The formula \emph{the interesting is
interesting} is therefore not a universal tautology, but a claim of possible
self-application.

This loop partly explains why the concept is elusive. Every acquired grip
changes the repertoire of examples, capacities for discrimination, and frontier
of the problem; a productive definition thus brings forth the cases that will
require its revision. Neither paradox nor ineffability follows. Local and
controllable properties can still be formulated, but their theory must be
reflexive, historical, and revisable.

Finally, self-application is not a criterion of truth. An interesting theory of
the interesting may be false; a boring theory may be accurate. Recursion
describes what the inquiry does to its object and its subject, not what validates
its propositions.

\section{Theoretical Framework and Programme of Naturalisation}

The theoretical framework will articulate three requirements. The first comes
from a philosophy of the subject: experience depends on conditions of reception
and judgement without being reducible to an objective feature. The second comes
from a phenomenology of temporality: retention, protention, bodily memory, and
aspect change describe the temporal constitution of a form. The third comes
from morphodynamic theories and complex systems: a global organisation has
properties that cannot be read from its isolated elements
\parencite{petitot1993phenomenologie,petitot2002naturaliser}.

The preface to Part III of the \emph{Ethics} provides a precedent for this
programme of naturalisation \parencite{spinoza1966ethique}. Spinoza refuses to
treat human affects as an ``empire within an empire'' exempt from the common
laws of nature. Studying the interesting takes up this gesture: a subjective
and situated experience can be explained through its causes without being
reduced to an intrinsic property of the object. Objectifying the phenomenon,
which makes its conditions and trajectories comparable, must be distinguished
from an objectivism that would locate interest in the form alone.

Ideas of affections involve both ``the nature of our body and the present nature
of the external body''
\parencite[III, prop. XXVII, dem.]{spinoza1861oeuvres}. This schema makes it
possible to think about a causality of encounter: the structure of the form, the
history and present state of the subject, the collective horizon, and the
circumstances jointly produce interest. The latter can be described as an
affective-cognitive regime of transition that changes the subject's power to
perceive, understand, anticipate, or act.

Naturalising this relation requires coordination between first-person
descriptions, behaviours, learning processes, formal structures, and
computational models, without assuming that any one level will exhaust the
others. The pleasures of discovery and construction
\parencite{feynman1999pleasure,florman1996existential} may accompany the
interesting without defining it: painful understanding may remain interesting,
whereas a pleasing explanation with no discriminating power may be only an
impression of clarity.

\part{Putting the Concept to Work}

\section{Music: Temporality and Micro-emotions}

Music is the principal laboratory of the inquiry because it makes relations
between memory, expectation, surprise, resolution, repetition, and exhaustion
especially visible. Meyer and later Narmour showed that musical affect and
meaning depend on learned expectations that are fulfilled, delayed, or thwarted
\parencite{meyer1956emotion,narmour1990basic,narmour1992complexity}.
Information-dynamic and neuro-symbolic models specify some of these processes
\parencite{abdallah2009information,berger1999expectations,gang1999unified}.

Explaining why a track is interesting cannot, however, amount to locating an
average quantity of surprise. The inquiry will examine trajectories of
attention at several scales: local micro-emotions, relations between melody and
harmony, global form, listening history, and the listener's own activity.
Jotney songs will provide a guiding case: they may be heard as solutions to the
problem of sustaining an autonomous, singable melody while moving the harmony
without displaying the procedure.

\emph{Histoire d'une oreille} provides a phenomenological vocabulary and a
method \parencite{pachet2018oreille}. \emph{Superimposed listening} submits
music to internal variations. An \emph{optale} is the feeling that a form is a
local optimum because nearby variants fail; a \emph{pseudoptale} warns that
familiarity may simulate this necessity. \emph{False evidence} and
\emph{overattentive listening} make it possible to examine false positives
produced by listeners themselves. Music will thus test the temporality, path
dependence, assimilable gap, non-monotonicity, and self-extinction of interest.

The good crossword clue described by Perec provides a linguistic analogue of
the optale \parencite{perec1999motscroises}. Its answer appears obvious once it
has been found, however insoluble the problem seemed before: the clue--answer pair then
appears as a local optimum that a more explicit formulation or another answer
would weaken. The case is crucial---even ``cruciverbal''---because it makes the
transformation from a plurality of readings to retrospective necessity
especially brief and observable.

\section{Artificial Intelligence: Explication and Loss of the Target}

Artificial-intelligence systems constitute a second laboratory. Work on Flow
Machines, reflexive interactions, Markov generation under constraints, and
computational interestingness has produced objects, protocols, successes, and
failures
\parencite{pachet2006clefs,pachet2008machines,pachet2012virtuosity,pachet2021assisted,pachet2026markov}.
It makes explicit operations that are often left implicit: sampling,
constraint, evaluation, model construction, selection, and direction.

AI is also the privileged reflexive case of the hostility of formalisation. A
system may solve the formal problem while removing its interesting part.
Metrics of novelty, surprise, diversity, preference, or popularity are local
indicators; they must not be confused with the relational target
\(I(F,S\mid H,t)\). Relevant experimental success will therefore not merely be
the generation of objects judged interesting, but the capacity of a model to
predict or explain transformations of attention, learning, and construction.

\section{Creativity: Producing, Judging, and Reconstructing a Problem}

Creativity makes it possible to distinguish the production of novelty from the
production of the interesting. An output may be unprecedented yet arbitrary; a
solution may use familiar materials while reconfiguring the space of
possibilities. Some forms appear necessary after they emerge even though they
could not have been deduced beforehand. They reveal their own constraints and
create an impression of surface simplicity that conceals an improbable genesis
\parencite{pachet2026impossibilite,pachet2026virtuosite}.

An interesting object may thus be encountered before its problem is known. The
subject reconstructs backwards the constraints satisfied by the solution, the
alternatives it excludes, and the question it makes visible. This reconstruction
must not become ad hoc: it should support variations, predictions, or
comparisons that could have failed. Creativity will particularly test
representational productivity, structured incompleteness, transformation of the
space of possibilities, and generative compression.

\section{Philosophy: Fertility, Truth, and Self-application}

Philosophical texts will not serve as disciplinary decoration. They will
constrain the concept, reveal its presuppositions, and provide counterexamples.
The historical corpus will show how the phenomenon has been distributed among
judgement, aesthetics, curiosity, attention, invention, truth, and boredom.
Heinz Wismann's remark will in particular prevent us from confusing a text's
capacity to set thought in motion with the validity of its propositions.

Philosophy is also the field of self-application. A theory of the interesting
must explain why it itself triggers distinctions and renewals, while accepting
that this productivity does not validate it. Its evaluation will therefore
require two sets of criteria: conceptual yield, production of problems, and
transfers on one hand; coherence, historical accuracy, and resistance to
objections on the other.

\section{Literature: Incompleteness, Style, and Interpretation}

Literature provides a field distinct from music because construction is
mediated by language, narrative, voices, and semantic expectations. It will
allow us to study how a sentence or narrative supplies enough grips to orient
interpretation without exhausting it. Interest may arise from structured
incompleteness, narrative expectation, formal constraint, a ruse with language,
or a displacement of ordinary categories.

This field must also distinguish fertility from pseudo-depth. Opacity may
generate interpretations because it is structured; it may also sustain only the
indefinite promise that a key will eventually arrive. Stylistic variations,
paraphrases, competing continuations, and comparisons between readings will
help determine whether a text produces transferable grips or merely
fascination. Literature will thus test structured incompleteness, mediated
contagion, dependence on the cultural horizon, and interpretive recursion.

\section{Comparative Method and Criteria of Testing}

The five fields are not intended to illustrate after the event an already
completed theory. Each must constrain the concept, reveal a limit, or produce a
distinction. The research will combine:

\begin{itemize}
  \item conceptual analysis of the interesting and its neighbours;
  \item historical reconstruction and direct reading of sources;
  \item a situated phenomenology of musical, literary, mathematical,
        scientific, and philosophical cases;
  \item counterfactual comparisons: internal variations, merely competent
        continuations, repetitions, and failed imitations;
  \item reflexive analysis of constructed systems and, where relevant,
        experiments on repetition, abandonment, learning, and changing
        judgements;
  \item confrontation with models of expectation, compression, curiosity,
        flow, and constrained generation.
\end{itemize}

Three levels will remain distinct: scientific results and effective properties
of systems; philosophical propositions; and argued articulations explaining
what the construction of a system or analysis of a case makes thinkable without
claiming that it proves it.

A hypothesis will gain strength if it can describe a trajectory, identify the
conditions under which interest varies, distinguish phenomena previously
confused, anticipate a change, or produce a counterexample. The thesis will also
accept negative results: some properties may hold only in one field, and no
single model will be assumed capable of exhausting the object.

\section{Corpus}

The main corpus will bring together \emph{Histoire d'une oreille}, Jotney songs,
the essays on the impossibility of creation and on virtuosity
\parencite{pachet2018oreille,pachet2026impossibilite,pachet2026virtuosite}, work
on generative systems and computational interestingness, and the philosophical
and scientific sources actually used.

The corpus will be organised in three circles: sources directly argued with in
the manuscript; cases for comparison and objection; and a programme of readings
still to be consolidated. The literary corpus will be narrowed according to the
properties to be tested rather than expanded in pursuit of historical
exhaustiveness. This organisation must maintain an explicit link between each
source, the proposition it supports, and the actual status of its reading.

\section{Expected Contributions}

The thesis aims to make six contributions:

\begin{enumerate}
  \item a problematised history of the interesting and of its dispersion among
        competing concepts;
  \item a relational, dispositional, and constructive definition, expressed by
        \(I(F,S\mid H,t)\);
  \item an organised cluster of properties for comparing the appearance,
        maintenance, circulation, exhaustion, and renewal of interest;
  \item a reflexive theory explaining both the concept's self-application and
        part of its elusiveness;
  \item a series of differential tests in music, AI, creativity, philosophy,
        and literature;
  \item a method of local formalisation that explicitly distinguishes target,
        indicator, scientific result, and philosophical scope.
\end{enumerate}

The outcome will not necessarily be a closed definition. It may take the form
of an articulated system of properties, distinctions, and criteria of
variation. A formal extension will examine, in particular,
non-compositionality, non-monotonicity, dynamic modalities, and the conditions
of a minimal axiomatics.

\section{Current Thesis Plan}

\begin{enumerate}
  \item \textbf{Part One --- Establishing the problem}
    \begin{enumerate}
      \item The blind spot and competing concepts.
      \item Genealogy and state of the art of the interesting.
      \item The interesting and the true.
      \item The hostility of formalisation.
    \end{enumerate}
  \item \textbf{Part Two --- Constructing the concept}
    \begin{enumerate}
      \item The relation \(I(F,S\mid H,t)\) and the construction of a grip.
      \item Relational, dynamic, and transmissive properties.
      \item Learning, flow, the almost learnable, and language.
      \item Recursion and self-application: ``the interesting is interesting.''
      \item Naturalisation, limits, and local models.
    \end{enumerate}
  \item \textbf{Part Three --- Putting the concept to work}
    \begin{enumerate}
      \item Music: expectations, micro-emotions, and variations.
      \item Artificial intelligence: generation, evaluation, and loss of the
            target.
      \item Creativity: invented necessity and implicit problem.
      \item Philosophy: fertility, truth, and reflexivity.
      \item Literature: incompleteness, style, and interpretation.
      \item Comparison of the fields and revision of the concept.
    \end{enumerate}
\end{enumerate}

\section{Directly Related Personal Work}

\begin{itemize}
  \item \textcite{pachet2026virtuosite} analyses forms of virtuosity and their
        intended effects on listeners.
  \item \textcite{pachet2026impossibilite} examines scientific and conceptual
        obstacles to the idea of creating novelty.
  \item \textcite{pachet2026markov} synthesises work on Markov modelling and
        constraint-based controlled generation.
  \item \textcite{pachet2021assisted} draws new categories of novelty from the
        Flow Machines projects.
  \item \textcite{pachet2018oreille} describes the ontogenesis of a musical ear
        and provides the attentional concepts used here.
  \item \textcite{pachet2012virtuosity} connects virtuosity, creativity, and the
        construction of the Virtuoso system.
  \item \textcite{pachet2008machines} presents reflexive systems as machines
        favouring departure from oneself.
  \item \textcite{pachet2006clefs} offers an initial model of interesting
        melodies.
\end{itemize}

\printbibliography[title={References}]

\end{document}
